\section{Actual and fresh projections through independent kernels}
\label{rhyb:section}\label{sec:r-hybrids}

We compare the information in a block's incoming and ending states after
revealing independent auxiliary measurements of the same signal.  The
comparison averages over the number of auxiliary rows.

We use the uniform-sphere experiment in
Definition~\ref{def:r-streaming-model} and the entropy conventions of
Section~\ref{sec:r-experiments}.  Write \(\gamma=N(0,I_d)\) for
the law of a single row.  All intermediate lemmas concern the
fixed-seed experiment supplied by Lemma~\ref{lem:r-seed-reduction}
and the Borel reduction in Lemma~\ref{lem:r-borel-kernels}; in
particular initialization remains jointly independent of the signal
and all rows.  Transitions may use fresh randomness, and the terminal
output uses only its state and index.  Auxiliary rows are analysis
variables, independent of the joint signal and learner variables as
specified below.  Constants \(c,C>0\) are absolute and may increase.
The estimates are for sufficiently large integer \(d\).

For a nonnegative integer \(j\), write \(D'_j=(Z,ZS)\) for \(j\) fresh
measurements.  The matrix \(Z\) is independent of the joint vector consisting of the
signal and all learner variables already generated; \(D'_0\) is empty.  Set
\begin{equation}\label{rhyb:parameters}
 k_0=\lceil d/4\rceil,\qquad k_1=\lfloor d/2\rfloor,\qquad
 N_k=k_1-k_0+1,\qquad h_0=\lfloor d/8\rfloor,\qquad
 q=\lfloor d/32\rfloor.
\end{equation}
For a finite state \(V\), with fresh rows independent of \((S,V)\), define
\begin{equation}\label{rhyb:potential}
 J_k(V)=I(S;V\mid D'_k),\qquad
 \overline J(V)=\frac1{N_k}\sum_{k=k_0}^{k_1}J_k(V).
\end{equation}
These quantities are finite and nonnegative because they are at most
\(H(V)\).

Consider a block of \(1\le m\le h_0\) actual observations, with incoming
state \(U\) and ending state \(V\).  Assume \(H(V)\le B<\infty\).
The block rows \(X_1,\ldots,X_m\) are independent of \((S,U)\).  The
block transition produces \(V\) from \(U\) and the block data by a
measurable kernel, conditionally independently of \(S\).

\begin{proposition}[Averaged information increase per block]
\label{rhyb:block}
For the block above, with \(1\le m\le h_0\) and \(H(V)\le B\),
\begin{equation}\label{rhyb:block-cost}
 \overline J(V)-\overline J(U)
 \le Cm+\frac{m}{N_k}\left(C+\frac{B}{h_0}\right)+\frac{B}{q}.
\end{equation}
When \(B\le d^2\), the right side is at most \(C'd\) for an absolute
constant \(C'\).
\end{proposition}

Keeping \(V\) as produced by the actual block, replace fresh auxiliary
rows one at a time by rows the learner actually saw.  Conditioning on the
ending state can bias an actual row; the total row-information cost is at
most \(H(V)\).
Independent random kernels control the accompanying entropy cost, and
averaging over the auxiliary count reduces the difference of successive
conditional entropies to two endpoint terms.

\subsection{Residual geometry and finite label entropies}

The following residual-sphere estimates justify the entropy subtractions
and will also normalize the row comparison.

\begin{lemma}[Residual-sphere estimates]\label{rhyb:residual}
Let \(D_j=(X,XS)\) consist of \(j\) independent Gaussian measurements,
where \(0\le j<d\), and put \(n=d-j\).  Conditional on \(D_j\), the signal
is uniform on a sphere of radius \(R\) in the affine solution space.
Marginally,
\begin{equation}\label{rhyb:radius}
 R^2\stackrel{\mathrm{law}}=
 \frac{G_n}{G_n+G_j},
 \qquad
 \mathbb E R^{-2}=1+\frac{j}{n-2}\quad(n>2),
\end{equation}
where \(G_n,G_j\) are independent chi-squared variables and \(G_0=0\).

If \(j\le 3d/4\), the entropy of one further label given its row and
\(D_j\) lies between two absolute constants, and its conditional log
density is absolutely integrable.  Joint label entropies, and the
conditional label entropies obtained by also conditioning on a finite
state \(V\), are finite whenever the total number of rows is at most
\(3d/4\).  The associated conditional log densities are absolutely
integrable.  Removing \(V\) can increase any such joint entropy by at
most \(H(V)\).

Let \(\sigma_R\) be uniform probability on a sphere of radius \(R>0\) in an
\(n\)-dimensional affine space.  A Euclidean ball of radius \(a\) with
arbitrary center has \(\sigma_R\)-mass at most
\((2a/R)^{n-1}\) when \(2a/R\le 1/2\).  For \(n\ge3\),
\begin{equation}\label{rhyb:reciprocal-sphere}
 \sup_{s\in\operatorname{supp}\sigma_R}
 \int\frac{\sigma_R(dt)}{\|t-s\|}
 \le \frac{C}{R}.
\end{equation}
\end{lemma}

\begin{proof}
The rows have full rank almost surely.  Their labels determine the
row-space component of \(S\).  Rotational invariance makes the remaining
direction uniform in the kernel, conditionally on its length and the
data.  Representing a uniform sphere point as a standard Gaussian vector
divided by its norm proves the ratio law in \eqref{rhyb:radius}.
For \(n>2\), independence and the chi-squared density give
\(\mathbb E G_n^{-1}=1/(n-2)\) and \(\mathbb E G_j=j\), proving the inverse
moment.

Let \(P\) be projection onto the kernel and let \(x\) be a new row.
Conditional on \(D_j,x\), the new label is a fixed shift plus
\[
 R\|Px\|U_1,
\]
where \(U\) is uniform on \(S^{n-1}\).  The density of \(U_1\) is
proportional to \((1-u^2)^{(n-3)/2}\) on \((-1,1)\).  Its normalizing
integral is at least \(c/\sqrt n\), by integrating over
\(|u|\le1/\sqrt n\).  For the dimensions in use this density is at most
\(C\sqrt n\), and hence
\[
 h(U_1)\ge-\log(C\sqrt n).
\]
Its log density is integrable: near an endpoint this is the
integrability of \(t^\alpha|\log t|\) for \(\alpha>-1\).

Also \(\|Px\|^2\) is chi-squared with \(n\) degrees of freedom.
For \(j\le3d/4\), the inverse moments of \(R\) and \(\|Px\|/\sqrt n\)
are bounded by absolute constants.  Jensen's inequality gives, for
example,
\[
 \mathbb E\log R\ge-\tfrac12\log\mathbb E R^{-2}\ge-C,
 \qquad
 \mathbb E\log(\|Px\|/\sqrt n)\ge-C.
\]
Ordinary positive moments control the positive logarithmic parts.
The entropy formula under scaling now gives an absolute lower bound for
the conditional entropy of the next label.  The upper bound is
\(\tfrac12\log(2\pi e)\), the entropy of its marginal standard normal
law.  These calculations also give absolute integrability of the
conditional log density.  Applying the chain rule in any order gives
the same conclusion for all ordered subsets of labels with the stated
row budget.  Extra independent rows may be conditioned on before \(V\)
is added, since they do not change these baseline distributions.

Adding the finite state reduces a joint label entropy by a conditional
mutual information at most \(H(V)\).  To check absolute integrability as
well, let \(r=dQ/dP\) be the density ratio of the conditional joint law
to its product-of-conditionals reference law.  The negative part of
\(\int\log r\,dQ\) is at most \(1/e\), because
\(-r\log r\le1/e\) for \(0<r\le1\).  Finite relative entropy therefore
also makes the positive part integrable.  The log density after
conditioning on \(V\) is the baseline log density plus this integrable
log ratio.  All label-entropy chain rules just used are consequently
between finite quantities.

For \(n=1\) the cap bound is the trivial bound one.  For \(n\ge2\), first
scale to \(R=1\) and center the ball at a
sphere point.  If its radius is \(a\le1/2\), its projection onto the
tangent coordinate plane lies in an \((n-1)\)-ball of radius \(a\).
The graph area factor there is at most \(2\), while the whole sphere
has area at least twice the volume of the unit \((n-1)\)-ball.
The cap has mass at most \(a^{n-1}\).  For an arbitrary center, choose a
sphere point in the intersection, if one exists, and double the radius.
This proves the displayed arbitrary-center bound.

Finally, for a sphere point \(s\), integrate the centered cap estimate
against \(t^{-2}\,dt\) for \(0<t\le R/2\); the integral is bounded by
\[
 R^{-(n-1)}\int_0^{R/2}t^{n-3}\,dt\le C/R
\]
when \(n\ge3\).  The part with \(t>R/2\) is at most \(C/R\) using total
mass one.  The layer-cake formula for \(\|t-s\|^{-1}\) proves
\eqref{rhyb:reciprocal-sphere}.
\end{proof}

\subsection{A block reduction and its row-information cost}

For this block, write
\[
 X_{\le i}=(X_1,\ldots,X_i),\qquad
 Y_{\le i}=X_{\le i}S,\qquad
 D_i=(X_{\le i},Y_{\le i}),\qquad L_i=(V,D_i)
\]
for \(0\le i\le m\), with the evident empty tuples at \(i=0\).
The rows are Gaussian before conditioning on \(V\); their conditional
laws after \(V\) is known need not be Gaussian.

For an independent sequence of fresh rows \(z_1,z_2,\ldots\), define the
successive conditional entropies
\begin{equation}\label{rhyb:slopes}
 a_j^{(i)}
 =h(z_jS\mid z_j,D'_{j-1},L_i),
 \qquad j\ge1,\quad i+j\le3d/4.
\end{equation}
All the entropies used below have a total row count at most \(3d/4\), so
Lemma~\ref{rhyb:residual} applies.

\begin{lemma}[One-row hybrid identity]\label{rhyb:hybrid}
For \(k_0\le k\le k_1\) and \(0\le i\le m\), let \(Z^{(i,k)}\) be a fresh
matrix of \(k-i\) rows and define
\[
 G_i(k)=
 h\bigl((Y_{\le i},Z^{(i,k)}S)\mid X_{\le i},Z^{(i,k)},V\bigr)
 +I(S;X_{\le i}\mid V).
\]
For \(i<m\), put \(x=X_{i+1}\), \(Y=xS\), \(L=L_i\), and \(r=k-i-1\).
With \(F_r=(Z,ZS)\) denoting common fresh data of size \(r\),
\begin{equation}\label{rhyb:hybrid-step}
 G_i(k)-G_{i+1}(k)
 =a_{r+1}^{(i)}
 -\bigl[h(Y\mid x,F_r,L)+I(S;x\mid F_r,L)\bigr].
\end{equation}
The actual-row divergence costs are
\begin{align}
 K_i&:=\mathbb E D_{\rm KL}
 \bigl(\mathcal L(X_{i+1}\mid S,L_i)\,\|\,\gamma\bigr)
 =I(X_{i+1};V\mid S,X_{\le i}),\label{rhyb:row-cost}\\
 \sum_{i=0}^{m-1}K_i
 &=I(X_1,\ldots,X_m;V\mid S)\le H(V)\le B.\label{rhyb:row-budget}
\end{align}
Finally,
\begin{equation}\label{rhyb:block-reduction}
 J_k(V)\le J_k(U)+G_0(k)-G_m(k).
\end{equation}
\end{lemma}

\begin{proof}
We first relate \(G_i(k)\) to conditional information without
subtracting differential entropies of the singular signal.  Suppose
\(D=(A,AS)\) has the marginal law of \(k\) independent Gaussian
measurements, although its rows may have helped produce \(V\).  The
conditional relative-entropy chain rule gives
\[
 I(S;V\mid D)-I(S;V)=I(D;V\mid S)-I(D;V).
\]
Every term here is finite, being bounded by \(H(V)\).  The labels are
determined by \(A,S\), and \(A\) is marginally independent of \(S\).
Expanding the finite terms involving \(V\) therefore yields
\begin{equation}\label{rhyb:information-identity}
 I(S;V\mid D)=I(S;V)-h(AS\mid A)+h(AS\mid A,V)+I(S;A\mid V).
\end{equation}
The label entropies are finite by Lemma~\ref{rhyb:residual}.
In a hybrid, the fresh rows contribute zero to the last mutual
information.  All hybrids have the same baseline entropy \(h(AS\mid A)\).
Thus
\[
 J_k(V)-I(S;V\mid D_m,D'_{k-m})=G_0(k)-G_m(k).
\]
Given \(D_m,D'_{k-m}\), the block transition is a Markov kernel from
\(U\) to \(V\), conditionally independently of \(S\).  Conditional
data processing gives
\[
 I(S;V\mid D_m,D'_{k-m})
 \le I(S;U\mid D_m,D'_{k-m})=J_k(U).
\]
The equality holds because these \(k\) rows are fresh relative to
\((S,U)\).  This proves \eqref{rhyb:block-reduction}.

For the one-row identity, realize the two adjacent hybrids using the
same \(r\) fresh rows \(Z\), one further fresh row for \(G_i(k)\), and
the actual row \(x\) for \(G_{i+1}(k)\).  Expand the labels in the order
\(Y_{\le i},ZS,\text{last label}\).  Before the last row is added the
common part is
\[
 h(Y_{\le i}\mid X_{\le i},V)
 +I(S;X_{\le i}\mid V)+h(ZS\mid Z,L).
\]
The remaining term for \(G_i(k)\) is \(a_{r+1}^{(i)}\).
For \(G_{i+1}(k)\), conditioning the first two label groups on the
actual row subtracts, respectively,
\[
 I(Y_{\le i};x\mid X_{\le i},V)
 \quad\hbox{and}\quad I(ZS;x\mid Z,L).
\]
Its new row-information term is initially
\(I(S;x\mid X_{\le i},V)\).  The first subtraction leaves \(I(S;x\mid L)\),
because the old labels are determined by \(S,X_{\le i}\).
The second leaves \(I(S;x\mid F_r,L)\): the fresh matrix \(Z\) is
independent, and \(ZS\) is determined by \(S,Z\).  The last label then
contributes \(h(Y\mid x,F_r,L)\).  This proves \eqref{rhyb:hybrid-step}
and retains both earlier-label subtractions.

Conditional on \(S,X_{\le i}\), the unconditioned law of \(X_{i+1}\)
is \(\gamma\), and the labels in \(L_i\) are redundant.  The definition
of conditional mutual information proves \eqref{rhyb:row-cost}.
Summing it is the conditional chain rule for \(X_1,\ldots,X_m\).
The resulting information is at most \(H(V)\), proving
\eqref{rhyb:row-budget}.  In particular the conditional actual-row
laws have finite divergence from \(\gamma\) almost surely.
\end{proof}

The identity has isolated the issue caused by conditioning on \(V\).
We next lower-bound the bracket in \eqref{rhyb:hybrid-step}.  Its reference
posterior will exclude the actual row, and the change from a Gaussian
row to the actual row will be charged through \(K_i\).

\subsection{Independent kernels and the biased-row comparison}

Two signals giving the same fresh labels differ by a vector in the
fresh matrix's kernel.  The posterior will select such directions
nonuniformly, so we need an inverse-volume estimate that holds for every
selection rule allowed below.

\begin{lemma}[Inverse volumes for independent kernels]
\label{rhyb:inverse-volume}
Let \(r,q\) be integers satisfying
\[
 d/8\le r<d,\qquad 1\le q\le r/2,\qquad r-q\ge2.
\]
Generate independent standard Gaussian \(r\)-by-\(d\) matrices
\(Z_1,\ldots,Z_q\).  After generating \(Z_j\), select a unit vector
\(u_j\in\ker Z_j\) by a probability kernel that may depend on \(Z_j\)
and all preceding matrices and selections, but not on future matrices.
Then
\begin{equation}\label{rhyb:volume-moment}
 \mathbb E\,\operatorname{vol}(u_1,\ldots,u_q)^{-1}\le C^q,
\end{equation}
uniformly over all such selection kernels.  In particular the volume
is positive almost surely.
\end{lemma}

\begin{proof}
Fix a previous span \(P\) of dimension \(p<q\), and let \(E\) be the row
space of the next matrix.  When \(p>0\), set
\[
 \xi(P,E)=\inf_{\substack{v\in P\\\|v\|=1}}\|P_Ev\|,
\]
and set \(\xi=1\) for \(p=0\).  An operator and its adjoint have the same
norm, so
\(\|P_P|_{E^\perp}\|=\|P_{E^\perp}|_P\|\).  For \(p>0\), taking
complements of their squared norms shows, for every unit \(u\in E^\perp\),
that
\begin{equation}\label{rhyb:adjoint}
 \operatorname{dist}(u,P)^2
 \ge1-\|P_P|_{E^\perp}\|^2
 =\inf_{\substack{v\in P\\\|v\|=1}}\|P_Ev\|^2
 =\xi(P,E)^2.
\end{equation}
For \(p=0\), both the distance and \(\xi\) equal one, so all the bounds
below are immediate.  We may therefore assume \(p>0\) in the small-ball
calculation.

The new \(E\) is a uniform \(r\)-dimensional subspace independent of
the preceding selections.  For a fixed unit \(v\) and \(0<t\le1/4\),
\begin{equation}\label{rhyb:fixed-vector-small-ball}
 \mathbb P\{\|P_Ev\|\le2t\}\le(Ct)^r.
\end{equation}
Indeed \(\|P_Ev\|\) has the distribution of
\(\|g\|/(\|g\|^2+\|g'\|^2)^{1/2}\), with independent standard Gaussian
vectors of dimensions \(r\) and \(d-r\).  The event forces
\(\|g\|\le Ct\|g'\|\).  The Gaussian integral, restricted to the ball
of radius \(\sqrt r\), gives
\((2\pi)^{r/2}\ge e^{-r/2}v_r r^{r/2}\), hence
\(v_r\le(C/\sqrt r)^r\).  Conditional on \(g'\), this volume estimate
and the Gaussian density bound make its probability at most
\((Ct\|g'\|/\sqrt r)^r\).
For \(z\ge0\),
\[
 z^{r/2}\le(2r/e)^{r/2}e^{z/4},
 \qquad
 \mathbb E e^{\|g'\|^2/4}=2^{(d-r)/2}.
\]
Since \(r\ge d/8\), these estimates give
\(\mathbb E\|g'\|^r\le(C\sqrt d)^r\).  Averaging proves
\eqref{rhyb:fixed-vector-small-ball}.

A maximal \(t\)-separated subset of the unit sphere in \(P\) is a
\(t\)-net with at most \((1+2/t)^p\le(3/t)^p\) points, by packing
disjoint balls of radius \(t/2\) in the ball of radius \(1+t/2\).
Projection is a contraction, so \(\xi\le t\) puts some net point at
projection distance at most \(2t\).  By \eqref{rhyb:fixed-vector-small-ball},
\begin{equation}\label{rhyb:subspace-small-ball}
 \mathbb P\{\xi\le t\}\le(3/t)^p(Ct)^r
 \le(C_0t)^{r-p},\qquad 0<t\le1/4.
\end{equation}
Here \(C_0\) is absolute because \(p<q\le r/2\).
Choose a fixed \(t_0\le1/4\) such that \(C_0t_0\le1/2\).
Since \(r-p>1\), integration of \eqref{rhyb:subspace-small-ball} gives
\[
 \mathbb E\xi^{-1}
 =1+\int_0^1\mathbb P\{\xi<t\}\,\frac{dt}{t^2}
 \le C.
\]
For \(t\ge t_0\), the trivial probability bound suffices.
The estimate is uniform in the preceding span \(P\).

Gram--Schmidt expresses the volume as the product of distances from
successive preceding spans.  Conditional on all previous selections,
the \(p=0\) observation and \eqref{rhyb:adjoint} bound the reciprocal of
the next distance by
\(\xi^{-1}\), whatever vector the new kernel selects.  The new row space
is still independent of that history.  Iterated conditional expectation
therefore bounds the inverse-volume product by \(C^q\), and also shows
that none of the distances vanishes with positive probability.
\end{proof}

Return to a fixed \(i,k\) in Lemma~\ref{rhyb:hybrid}.  Thus
\(x=X_{i+1}\), \(L=(V,D_i)\), \(r=k-i-1\), and \(F_r=(Z,ZS)\) is fresh.
Define the reference posterior and its reciprocal-distance scale by
\begin{equation}\label{rhyb:reciprocal-scale}
 \mu=\mathcal L(S\mid L,F_r),
 \qquad
 A(s,L,F_r)=\int\frac{\mu(ds')}{\|s'-s\|}.
\end{equation}
The posterior \(\mu\) includes the ending state but does not condition on
the actual row \(x\).  This distinction is essential: we will compare the
actual conditional law of \(x\) with a Gaussian row.

\Needspace{6\baselineskip}
\begin{lemma}[A row biased by the ending state]\label{rhyb:row-comparison}
In the preceding experiment, \(A(S,L,F_r)\ge1/2\) almost surely and
\(\mathbb E A(S,L,F_r)<\infty\) for each fixed \(d\).  For almost every
Gaussian row \(x\), the projection of \(\mu\) by \(s\mapsto xs\) has a
Lebesgue density \(p_{\mu,x}\).  With \(q=\lfloor d/32\rfloor\),
\begin{align}
 h(xS\mid x,L,F_r)+I(S;x\mid L,F_r)
 &\ge-\mathbb E\log p_{\mu,x}(xS),\label{rhyb:cross-entropy}\\
 \mathbb E\log\frac{p_{\mu,x}(xS)}{A(S,L,F_r)}
 &\le C+\frac{K_i}{q}.\label{rhyb:row-log-bound}
\end{align}
Both expectations use the actual row law.  The logarithms in these
displays are absolutely integrable.
\end{lemma}

\begin{proof}
Conditional on \(D_i,F_r\), before \(V\) is imposed, the signal is uniform
on a residual sphere of affine dimension
\[
 n=d-i-r=d-k+1>2.
\]
For a state \(v\) with
\(\pi_v=\mathbb P\{V=v\mid D_i,F_r\}>0\), its posterior density relative
to the uniform residual law is at most \(1/\pi_v\).  This follows directly
from Bayes' formula, since the conditional probability of \(V=v\) given
the signal and these data is at most one.  Lemma~\ref{rhyb:residual}
therefore gives
\[
 A(s,L,F_r)\le\frac{C}{R\pi_v}
\]
for every point \(s\) on the residual sphere.  Average first over \(s\)
and \(v\).  The factor \(\pi_v\) cancels, and the sum has at most \(W\)
terms if \(V\) has at most \(W\) values.  Consequently
\begin{equation}\label{rhyb:scale-integrability}
 \mathbb E A(S,L,F_r)\le CW\,\mathbb E R^{-1}<\infty.
\end{equation}
The inverse moment in \eqref{rhyb:radius} proves the last inequality.
Every pairwise distance on the unit sphere is at most two, giving
\(A\ge1/2\).  In particular \(\log A\) is absolutely integrable.

The posterior is nonatomic and is absolutely continuous relative to the
uniform residual law.  For every row whose projection into the residual
kernel is nonzero, the projection of that uniform law has a Lebesgue
density, by the one-coordinate formula in Lemma~\ref{rhyb:residual}.
The exceptional rows have Gaussian measure zero.  Absolute continuity
passes to the projected posterior, proving the assertion about
\(p_{\mu,x}\).

The information cost of conditioning this posterior on \(x\) is finite.
Indeed \(x\) and \(F_r\) are conditionally
independent given \(S,L\), so the conditional chain rule yields
\[
 I(S;x\mid L)-I(S;x\mid L,F_r)=I(x;F_r\mid L)\ge0.
\]
Also the chain rule for the relative entropy in \(K_i\) gives
\[
 K_i=I(S;x\mid L)+
 \mathbb E D_{\rm KL}\bigl(\mathcal L(x\mid L)\,\|\,\gamma\bigr).
\]
Thus
\begin{equation}\label{rhyb:information-charge}
 I(S;x\mid L,F_r)\le I(S;x\mid L)\le K_i.
\end{equation}
The conditional actual law of \(S\) given \(x,L,F_r\) is therefore dominated
by \(\mu\) almost surely.  Projecting by \(s\mapsto xs\) contracts its
relative entropy, so the actual projected law is dominated by the
reference projected law with expected divergence at most
\(I(S;x\mid L,F_r)\).  Lemma~\ref{rhyb:residual} supplies an integrable
log density for the actual conditional label.  Applying the
log-density-ratio argument from that lemma to the finite projected
divergence shows that \(\log p_{\mu,x}(xS)\) is absolutely integrable
as well.  We may now expand the projected divergence:
\[
 -\mathbb E\log p_{\mu,x}(xS)-h(xS\mid x,L,F_r)
 \le I(S;x\mid L,F_r).
\]
This is \eqref{rhyb:cross-entropy}.

It remains to prove \eqref{rhyb:row-log-bound}.  Fix a typical pair
\((s,L)\) under the joint law of \((S,L)\).  Generate \(q\) independent
copies \(Z_1,\ldots,Z_q\) of the \(r\) fresh rows, and put
\(F^{(j)}=(Z_j,Z_js)\).  Let \(\mu_j,A_j\) be the corresponding
posterior and scale from \eqref{rhyb:reciprocal-scale}.  Initially let
\(x\sim\gamma\) be independent of all these matrices.  We claim
\begin{equation}\label{rhyb:product-moment}
 \mathbb E_{Z_1,\ldots,Z_q,x}
 \prod_{j=1}^q\frac{p_{\mu_j,x}(xs)}{A_j}\le C^q.
\end{equation}

For \(\delta>0\), replace the density in the \(j\)-th factor by its
symmetric slab average
\[
 p^\delta_{\mu_j,x}(xs)
 =\frac{1}{2\delta}
 \mu_j\{s':|x(s'-s)|\le\delta\}.
\]
Draw \(s'_j\) separately from \(\mu_j\), and write
\(\Delta_j=s'_j-s\), \(u_j=\Delta_j/\|\Delta_j\|\).  The posteriors are
nonatomic, so the differences are nonzero almost surely.  Conditional on
these differences, the vector
\((x\Delta_1,\ldots,x\Delta_q)\) is Gaussian.  Its density is bounded by
\begin{equation}\label{rhyb:gaussian-density}
 (2\pi)^{-q/2}
 \prod_{j=1}^q\|\Delta_j\|^{-1}
 \operatorname{vol}(u_1,\ldots,u_q)^{-1};
\end{equation}
we may use \(+\infty\) for a singular tuple.  This is the usual Gaussian
density bound with covariance matrix
\((\langle\Delta_j,\Delta_\ell\rangle)_{j,\ell}\).
Its probability in \([-\delta,\delta]^q\) is at most
\eqref{rhyb:gaussian-density} times \((2\delta)^q\).

After division by \(\prod_j A_j\), nonnegative integration converts the
resulting bound into \((2\pi)^{-q/2}\) times an inverse-volume expectation.
Conditional on each \(Z_j\), the selected point in that expectation has
the probability law
\begin{equation}\label{rhyb:normalized-kernel}
 \nu_j(ds')=\frac{\mu_j(ds')}{\|s'-s\|A_j}.
\end{equation}
Each law is normalized by its own \(A_j\).  Its total mass is one for
every typical \(Z_j\), so the marginal law of the matrices remains the
product of their original Gaussian laws.  In particular the matrices
have not been tilted.  Both \(s\) and \(s'_j\) satisfy the labels \(Z_js\);
therefore \(u_j\in\ker Z_j\).  These selections, made separately
conditional on their matrices and \((s,L)\), satisfy the hypotheses of
Lemma~\ref{rhyb:inverse-volume}.  Indeed
\[
 r=k-i-1\ge k_0-h_0\ge d/8,\qquad q\le r/2,\qquad r-q\ge2
\]
for sufficiently large \(d\).  The lemma bounds the smoothed product
expectation by \(C^q\), uniformly in \(\delta\).

Here is a density version that makes passage to exact evaluations
legitimate.  In the unbiased one-copy experiment, \(x\) is independent
of \((S,L,F_r)\), so conditional on \(L,F_r,x\), the signal has law
\(\mu\).  The one-dimensional Lebesgue differentiation theorem applied
to the integrable density \(p_{\mu,x}\) shows that \(xS\) is almost surely
a Lebesgue point and that the symmetric slab averages converge to
\(p_{\mu,x}(xS)\).  Each slab average is jointly measurable in its
kernel parameters and evaluation point.  Taking its limit along
\(\delta=2^{-n}\), and assigning zero when no finite limit exists,
therefore gives a jointly measurable density version that agrees at
these evaluations.  Fubini's theorem gives the same assertion for
almost every fixed \((s,L)\).  For such a pair it holds simultaneously
for all \(q\) copies, since \(q\) is finite.  Fatou's lemma now proves
\eqref{rhyb:product-moment}.

Change only \(x\) from its unbiased law to its actual conditional law
given \(S=s,L\).  The \(Z_j\) remain independent of it and of one
another.  The divergence from the unbiased product experiment is
exactly
\(D_{\rm KL}(\mathcal L(x\mid s,L)\,\|\,\gamma)\).
For probability laws \(Q\ll P\) and a nonnegative function \(g\), the
entropy variational inequality gives
\begin{equation}\label{rhyb:entropy-inequality}
 \mathbb E_Q\log g
 \le D_{\rm KL}(Q\|P)+\log\mathbb E_P g.
\end{equation}
For positive bounded \(g\) this follows from Jensen's inequality applied
to \(g/(dQ/dP)\); truncation gives the extended form.  This is the
standard variational identity recorded in
\cite[Equations~(1.1)--(1.2)]{DupuisMao2022}.
Apply it to the product in \eqref{rhyb:product-moment}.
Finite row divergence transfers the unbiased almost-sure density
statements to the actual law.  The logarithms are integrable by
\eqref{rhyb:cross-entropy} and \eqref{rhyb:scale-integrability}.
The expectation of the log product is \(q\) times the one-copy
expectation, since the \(q\) fresh copies have the same marginal law.
Divide by \(q\) and average over \((S,L)\).  Equation
\eqref{rhyb:row-cost} then gives \eqref{rhyb:row-log-bound}.
\end{proof}

The independent-kernel estimate bounds the actual projected density at
the scale \(A\).  To compare this scale to a fresh conditional entropy,
we use two fresh rows after the \(r\) common rows.  The second of the two
is the next slope beyond the one already present in the hybrid identity.

\subsection{A reciprocal-distance tilt and two fresh rows}

\begin{lemma}[The next fresh entropy slope]\label{rhyb:tilted-coupling}
For the posterior and scale in \eqref{rhyb:reciprocal-scale},
\begin{equation}\label{rhyb:next-slope}
 a_{r+2}^{(i)}\le-\mathbb E\log A(S,L,F_r)+C.
\end{equation}
\end{lemma}

\begin{proof}
Write \(C_*=(L,F_r)\).  Conditional on \(C_*\), the signal has law
\(\mu\).  For an integer \(N\ge1\), define
\[
 w_N(s,t)=\min\{N,\|t-s\|^{-1}\},\qquad
 A_N(s,C_*)=\int w_N(s,t)\,\mu(dt),
\]
where \(w_N(s,s)=N\).  Given \(S=s,C_*\), draw an auxiliary point \(S^*\)
from the probability kernel
\[
 \nu_s^N(dt)=\frac{w_N(s,t)}{A_N(s,C_*)}\,\mu(dt).
\]
Distances on the unit sphere are at most two, so
\(1/2\le A_N\le N\) and \(d\nu_s^N/d\mu\le2N\).
The marginal law of \(S^*\) need not equal \(\mu\).

Compare the joint conditional law of \((S,S^*)\) with
\(\mu\otimes\mu\).  Its divergence is
\[
 \mathbb E\log\frac{w_N(S,S^*)}{A_N(S,C_*)}.
\]
The relative-entropy chain rule splits this into
\(I(S;S^*\mid C_*)\) plus the divergence of the
actual \(S^*\) marginal from the comparator \(\mu\).  Hence
\begin{align}
 I(S;S^*\mid C_*)
 &\le\mathbb E\log\frac{w_N(S,S^*)}{A_N(S,C_*)}\notag\\
 &\le-\mathbb E\log\|S-S^*\|
       -\mathbb E\log A_N(S,C_*).\label{rhyb:tilt-information}
\end{align}
The comparator divergence is finite, at most \(\log(2N)\).
The points differ almost surely because \(\mu\) is nonatomic.
Their log distance is absolutely integrable: its positive part is at
most \(\log2\); under \(\mu\otimes\mu\) its negative part is bounded by
an integrable multiple of \(\|S-S^*\|^{-1}\) by
\eqref{rhyb:scale-integrability}, and the tilted density is at most \(2N\).
Thus both inequalities in \eqref{rhyb:tilt-information} are between
well-defined finite quantities.

After drawing \(S^*\), generate independent fresh Gaussian rows
\(x_1,x_2\) and labels \(b_1=x_1S\), \(b_2=x_2S\).
Before \(S^*\) is added to the conditioning,
Lemma~\ref{rhyb:residual} gives a joint conditional density with an
integrable logarithm for these labels, since
\(i+r+2=k+1\le3d/4\).  Fresh-row independence and data processing give
\begin{align}
 I(b_2;S^*\mid b_1,x_1,x_2,C_*)
 &\le I((b_1,b_2);S^*\mid x_1,x_2,C_*)\notag\\
 &\le I(S;S^*\mid C_*).\label{rhyb:remove-coupling}
\end{align}
This is finite.  The log-density-ratio argument from
Lemma~\ref{rhyb:residual} consequently gives integrable conditional log
densities after \(S^*\) is added, both for the joint labels and for the
first label.  Their entropy chain rule justifies the conditional entropy
of the second label used next.

Given \(b_1,x_1,x_2,C_*,S^*\), compare the density of \(b_2\) with the Cauchy
density centered at \(x_2S^*\) and with scale
\(|b_1-x_1S^*|\).  The scale is positive almost surely.  With
\(\Delta=S-S^*\), the expected negative log of this reference density
at \(b_2\) is
\begin{align}
 &\log\pi+\mathbb E\log|x_1\Delta|
 +\mathbb E\log\left(1+\frac{(x_2\Delta)^2}{(x_1\Delta)^2}\right)
 \notag\\
 &\hspace{35mm}=\mathbb E\log\|\Delta\|+C_1,\label{rhyb:cauchy}
\end{align}
where \(C_1\) is absolute.  Conditional on nonzero \(\Delta\), the two
projections divided by \(\|\Delta\|\) are independent standard normals
\(G_1,G_2\).  The first logarithm is integrable because
\(\mathbb E|\log|G_1||<\infty\).  For the ratio term,
\[
 \log\bigl(1+(G_2/G_1)^2\bigr)
 =\log(G_1^2+G_2^2)-\log G_1^2,
\]
and both logarithms on the right are integrable by the one- and
two-dimensional chi-squared densities.  The cross-entropy inequality
therefore gives
\[
 h(b_2\mid b_1,x_1,x_2,C_*,S^*)
 \le\mathbb E\log\|\Delta\|+C_1.
\]

Removing \(S^*\) from the conditioning in this entropy bound and using
\eqref{rhyb:tilt-information}--\eqref{rhyb:remove-coupling} cancels
\(\mathbb E\log\|\Delta\|\).  We obtain
\[
 h(b_2\mid b_1,x_1,x_2,C_*)\le-\mathbb E\log A_N(S,C_*)+C_1.
\]
The left side is exactly \(a_{r+2}^{(i)}\) and is independent of \(N\).
Finally \(A_N\uparrow A\) and \(\log A_N\ge-\log2\).
Monotone convergence, with the integrability of \(\log A\) already
proved, yields \eqref{rhyb:next-slope}.
\end{proof}

\subsection{Averaging over the number of side measurements}

The row comparison and the two-row coupling now leave only a difference
of consecutive fresh entropy slopes.  Averaging over consecutive values
of \(k\) makes this difference telescope.

\begin{proof}[Proof of Proposition~\ref{rhyb:block}]
Equations~\eqref{rhyb:cross-entropy} and \eqref{rhyb:row-log-bound}
lower-bound the bracket in \eqref{rhyb:hybrid-step} by
\(-\mathbb E\log A-C-K_i/q\).  Equation~\eqref{rhyb:next-slope} then
gives
\[
 h(xS\mid x,L,F_r)+I(S;x\mid L,F_r)
 \ge a_{r+2}^{(i)}-C-K_i/q.
\]
Since \(r=k-i-1\), the hybrid identity becomes
\begin{equation}\label{rhyb:slope-difference}
 G_i(k)-G_{i+1}(k)
 \le a_{k-i}^{(i)}-a_{k-i+1}^{(i)}+C+K_i/q.
\end{equation}

For fixed \(i\), the fresh observations are exchangeable conditional on
\(L_i\).  Relabel two consecutive fresh observations and then condition
on the additional earlier one.  Conditional entropy can only decrease,
so the successive slopes \(a_j^{(i)}\) are nonincreasing in their
finite range.  Also
\[
 a_j^{(i)}\le\tfrac12\log(2\pi e)
\]
by comparison with the marginal standard normal label.

For \(i+j\le k_1+1\), a useful lower bound is
\begin{equation}\label{rhyb:slope-lower}
 a_j^{(i)}\ge-C-B/h_0.
\end{equation}
To prove it, use monotonicity to bound \(a_j^{(i)}\) below by the
average of the next \(h_0\) slopes.  Their sum is the joint entropy of
the next \(h_0\) fresh labels given their rows, \(D'_j\), and \(L_i\).
Removing \(V\) from that conditioning increases the entropy by at most
\(H(V)\le B\).  Without \(V\), the actual prefix and all fresh rows have
the ordinary independent-row law.  Each successive baseline label
entropy is at least \(-C\) by Lemma~\ref{rhyb:residual}, since the total
number of rows is at most
\[
 i+j+h_0\le k_1+1+h_0\le3d/4
\]
for sufficiently large \(d\).  The sum is therefore at least
\(-Ch_0-B\), proving \eqref{rhyb:slope-lower}.

Average \eqref{rhyb:slope-difference} over \(k_0\le k\le k_1\).
For each \(i\), its slope term is exactly
\[
 \frac1{N_k}\sum_{k=k_0}^{k_1}
 \bigl(a_{k-i}^{(i)}-a_{k-i+1}^{(i)}\bigr)
 =\frac{a_{k_0-i}^{(i)}-a_{k_1-i+1}^{(i)}}{N_k}
 \le\frac{C+B/h_0}{N_k}.
\]
Now sum over \(i=0,\ldots,m-1\), use
\(\sum_iK_i\le B\) from \eqref{rhyb:row-budget}, and apply
\eqref{rhyb:block-reduction}.  This gives \eqref{rhyb:block-cost}.
Finally \(N_k,h_0,q\ge cd\) and \(m\le d/8\), so \(B\le d^2\) makes
the displayed block cost at most \(C'd\).
\end{proof}


Subsection~\ref{subsec:r-hybrid-application} combines this block estimate with the stopping clock and the terminal residual-sphere bound.
